\documentclass[a4paper,12pt]{article}
\usepackage[T2A]{fontenc}
\usepackage[utf8]{inputenc}
\usepackage[russian]{babel}
\usepackage{amsmath,amssymb,mathtools}
\usepackage{PTSans}
\renewcommand{\familydefault}{\sfdefault}
\usepackage{icomma}
\usepackage[a4paper,margin=20mm,headsep=7mm,footskip=12mm]{geometry}
\usepackage{microtype}
\usepackage{tikz}
\usetikzlibrary{calc,arrows.meta,positioning,shapes.geometric}
\usepackage{tabularx,booktabs,enumitem}
\usepackage{hyperref}
\usepackage{parskip}
\usepackage{fancyhdr}
\usepackage{setspace}
\usepackage{xcolor}

\definecolor{accent}{HTML}{2F6F73}
\definecolor{darkaccent}{HTML}{19484B}
\definecolor{textgray}{HTML}{2E3336}
\definecolor{softgray}{HTML}{687579}
\definecolor{rulegray}{HTML}{CAD5D5}
\hypersetup{colorlinks=true,linkcolor=darkaccent,urlcolor=darkaccent}
\color{textgray}
\setstretch{1.32}
\setlength{\parindent}{0pt}
\setlist[itemize]{leftmargin=6mm,itemsep=0.8mm,topsep=1mm}
\setlist[enumerate]{leftmargin=7mm,itemsep=2.3mm,topsep=2mm}
\setlength{\headheight}{15pt}
\setlength{\tabcolsep}{5pt}
\renewcommand{\arraystretch}{1.25}
\pagestyle{fancy}
\fancyhf{}
\fancyhead[L]{\small\color{softgray}Смеси, растворы и проценты}
\fancyhead[R]{\small\color{softgray}11.10.26}
\fancyfoot[C]{\small\color{softgray}\thepage}
\renewcommand{\headrulewidth}{0pt}
\newcommand{\sectiontitle}[1]{%
  \vspace{1.5mm}{\Large\bfseries\color{darkaccent}#1}\par
  \vspace{0.8mm}{\color{rulegray}\hrule}\vspace{2.3mm}}
\newcommand{\key}[1]{\textbf{\color{darkaccent}#1}}
\newcommand{\formula}[1]{\[\displaystyle #1\]}
\newcommand{\unit}[1]{\,\text{#1}}

\begin{document}
\thispagestyle{empty}
\begin{center}
\vspace*{31mm}
{\small\color{softgray}ЧЕК-ЛИСТ}\par\vspace{7mm}
{\Huge\bfseries\color{darkaccent}Смеси, растворы\\[2mm]и проценты}\par
\vspace{7mm}
{\Large Сохранение вещества, таблицы, изменение цены}\par
\vspace{18mm}
{\large 11.10.26}\par
\vfill
{\large\bfseries Лёвин Артём Александрович}\par\vspace{2mm}
{\large эксклюзивно для Матвея Горбачёва}\par
\vspace{20mm}
\end{center}
\begin{tikzpicture}[remember picture,overlay]
\draw[accent,line width=1.2pt]
  ($(current page.south west)+(15mm,18mm)$)
  .. controls ($(current page.south west)+(68mm,32mm)$) and
  ($(current page.south east)+(-71mm,7mm)$) ..
  ($(current page.south east)+(-15mm,20mm)$);
\draw[rulegray,line width=0.7pt]
  ($(current page.south west)+(15mm,13mm)$)
  .. controls ($(current page.south west)+(71mm,25mm)$) and
  ($(current page.south east)+(-67mm,3mm)$) ..
  ($(current page.south east)+(-15mm,15mm)$);
\end{tikzpicture}
\clearpage

\sectiontitle{1. Одна модель для растворов, сплавов и смесей}

\key{Основная формула:} масса выбранного вещества равна массе всей смеси, умноженной на массовую долю этого вещества:
\formula{q=m\,c,\qquad c=\frac{p}{100},}
где $m$ --- масса смеси, $p$ --- содержание вещества в процентах, $c$ --- его доля (от $0$ до $1$), $q$ --- масса выбранного вещества.

\begin{center}
\begin{tabularx}{\linewidth}{@{}lXXX@{}}
\toprule
Объект & Масса смеси & Доля вещества & Масса вещества \\
\midrule
Первый сплав & $m_1$ & $c_1$ & $m_1c_1$ \\
Второй сплав & $m_2$ & $c_2$ & $m_2c_2$ \\
Полученный сплав & $m_1+m_2$ & $c$ & $(m_1+m_2)c$ \\
\bottomrule
\end{tabularx}
\end{center}

Если смешивают вещества без потерь, складываются как полные массы, так и массы выбранного компонента:
\formula{m=m_1+m_2,\qquad m\,c=m_1c_1+m_2c_2.}

\key{Пример.} Первый сплав содержит $10\%$ меди, второй --- $40\%$. Масса второго на $3\unit{кг}$ больше массы первого. После смешивания получен сплав с $30\%$ меди. Найдите массу полученного сплава.

\begin{center}
\begin{tabularx}{\linewidth}{@{}lXXX@{}}
\toprule
Сплав & Масса, кг & Доля меди & Медь, кг \\
\midrule
Первый & $x$ & $0{,}10$ & $0{,}10x$ \\
Второй & $x+3$ & $0{,}40$ & $0{,}40(x+3)$ \\
После смешивания & $2x+3$ & $0{,}30$ & $0{,}30(2x+3)$ \\
\bottomrule
\end{tabularx}
\end{center}

\formula{0{,}10x+0{,}40(x+3)=0{,}30(2x+3),\qquad x=3.}
\key{Ответ:} $2x+3=9\unit{кг}$. Вопрос относился ко \emph{всему} полученному сплаву, поэтому найденное $x$ ещё не является ответом.

\sectiontitle{2. Сушка: сохраняется сухое вещество}

При высушивании испаряется вода, а масса сухого вещества по условию остаётся постоянной. Долю сухого вещества находят вычитанием доли воды из единицы:
\formula{c_{\text{сух}}=1-c_{\text{воды}}.}

\key{Пример.} В винограде $90\%$ воды, в изюме $5\%$ воды. Сколько винограда необходимо для получения $20\unit{кг}$ изюма?
\formula{x\cdot 0{,}10=20\cdot 0{,}95,\qquad x=190\unit{кг}.}
\key{Идея:} состав меняется, но сухое вещество остаётся тем же. Следовательно, при сушке нужно приравнивать массы \emph{сухого вещества}, а не массы воды.

\sectiontitle{3. Две ситуации --- система двух уравнений}

Если условия описывают два эксперимента с одними исходными растворами, удобно составить \emph{по одной таблице для каждого опыта}. Неизвестные исходные массы в двух опытах одинаковы.

\key{Пример.} Смешали $40\%$ и $70\%$ растворы кислоты. Добавив $20\unit{кг}$ воды, получили $41\%$ раствор. Если вместо воды добавить $20\unit{кг}$ раствора кислоты с концентрацией $60\%$, получится $53\%$ раствор. Найдите массу первого раствора.

Пусть $x$ --- масса $40\%$ раствора, $y$ --- масса $70\%$ раствора.

\begin{center}
\begin{tabularx}{\linewidth}{@{}lXXX@{}}
\toprule
Первый опыт & Масса, кг & Доля кислоты & Кислота, кг \\
\midrule
Раствор 1 & $x$ & $0{,}40$ & $0{,}40x$ \\
Раствор 2 & $y$ & $0{,}70$ & $0{,}70y$ \\
Вода & $20$ & $0$ & $0$ \\
Итог & $x+y+20$ & $0{,}41$ & $0{,}41(x+y+20)$ \\
\bottomrule
\end{tabularx}
\end{center}

\begin{center}
\begin{tabularx}{\linewidth}{@{}lXXX@{}}
\toprule
Второй опыт & Масса, кг & Доля кислоты & Кислота, кг \\
\midrule
Раствор 1 & $x$ & $0{,}40$ & $0{,}40x$ \\
Раствор 2 & $y$ & $0{,}70$ & $0{,}70y$ \\
Добавленный раствор & $20$ & $0{,}60$ & $12$ \\
Итог & $x+y+20$ & $0{,}53$ & $0{,}53(x+y+20)$ \\
\bottomrule
\end{tabularx}
\end{center}

\formula{\begin{cases}
0{,}40x+0{,}70y=0{,}41(x+y+20),\\
0{,}40x+0{,}70y+12=0{,}53(x+y+20).
\end{cases}}

Вычтем первое уравнение из второго:
\formula{12=0{,}12(x+y+20)\quad\Rightarrow\quad x+y=80.}
Из первого уравнения $0{,}40x+0{,}70y=41$. Подставляя $x=80-y$, получаем $y=30$, $x=50$.

\key{Ответ:} $50\unit{кг}$ $40\%$ раствора.

\key{Важно.} Вода содержит $0\%$ кислоты. Если добавили раствор, его вклад в массу кислоты нужно учесть отдельно.

\sectiontitle{4. Когда неизвестна сама концентрация}

Иногда исходные массы одинаковы, а неизвестны доли вещества $a$ и $b$. Тогда общую массу можно выбрать условно: например, по $1\unit{кг}$ каждого раствора. При смешивании равных масс:
\formula{\frac{a+b}{2}=c_{\text{смеси}}.}
Если взять одну часть первого и три части второго, получим:
\formula{\frac{a+3b}{4}=c_{\text{смеси}}.}
Для ответа «сколько килограммов вещества?» потребуется ещё раз применить $q=mc$.

\sectiontitle{5. Изменение цены в процентах}

\key{Увеличение на $p\%$:} к исходной сумме $S_0$ прибавляют $p\%$ именно от неё. \key{Уменьшение на $p\%$:} из неё вычитают эту долю.
\formula{S_1=S_0\left(1+\frac{p}{100}\right)\quad\text{или}\quad
S_1=S_0\left(1-\frac{p}{100}\right).}

\begin{center}
\begin{tabularx}{\linewidth}{@{}lXXX@{}}
\toprule
Действие & Начальная цена & Изменение & Новая цена \\
\midrule
Рост на $11\%$ & $x$ & $+0{,}11x$ & $1{,}11x$ \\
Скидка $14\%$ & $2700$ & $-0{,}14\cdot2700$ & $0{,}86\cdot2700$ \\
\bottomrule
\end{tabularx}
\end{center}

\key{Пример 1.} Цена после повышения на $11\%$ составила $1443$ рубля. Начальная цена:
\formula{1{,}11x=1443\quad\Rightarrow\quad x=1300\unit{руб}.}

\key{Пример 2.} Цена снизилась с $2700$ до $2322$ рублей. Процент скидки считают \emph{от первоначальной цены}:
\formula{p=\frac{2700-2322}{2700}\cdot100\%=14\%.}

\key{Полезная проверка.} После скидки $4\%$ остаётся $96\%$ цены, то есть множитель равен $0{,}96$. Если цена уменьшилась в $4$ раза, скидка составила $75\%$.

\sectiontitle{6. Отношение стоимостей вместо отдельных цен}

В задачах сравнения часто достаточно найти \emph{отношение цен}. За $100\%$ принимают стоимость того объекта, относительно которого проводится сравнение.

\key{Пример.} Четыре рубашки дешевле куртки на $8\%$. На сколько процентов пять таких рубашек дороже куртки?

Пусть $r$ --- цена рубашки, $k$ --- цена куртки. Тогда
\formula{4r=0{,}92k\quad\Rightarrow\quad\frac{r}{k}=\frac{0{,}92}{4}.}
\formula{\frac{5r}{k}\cdot100\%=\frac{5\cdot0{,}92}{4}\cdot100\%=115\%.}
\key{Ответ:} пять рубашек дороже куртки на $115\%-100\%=15\%$.

\key{Два частых недочёта:}
\begin{itemize}
\item Найти неизвестное $x$, не проверив, что спрашивается: масса всей смеси, разность масс или масса вещества.
\item Принять $115\%$ за ответ в вопросе «на сколько процентов дороже»; требовалось именно превышение над $100\%$.
\end{itemize}

\clearpage
\sectiontitle{Алгоритм: как выбрать математическую модель}
\vspace{4mm}
\begin{center}
\begin{tikzpicture}[
  >=Stealth, font=\small,
  start/.style={ellipse,draw=darkaccent,thick,align=center,minimum width=38mm,minimum height=11mm,fill=white},
  act/.style={rounded corners=2mm,draw=rulegray,thick,align=center,text width=52mm,minimum height=14mm,inner sep=3mm},
  choice/.style={diamond,aspect=2.1,draw=darkaccent,thick,align=center,text width=37mm,inner sep=1.5mm},
  line/.style={-Stealth,thick,draw=accent}
]
\node[start] (start) at (0,0) {Прочитайте условие};
\node[act] (target) at (0,-1.8) {Определите, что известно\\и что нужно найти};
\node[choice] (decision) at (0,-4.8) {Какая величина\\меняется?};
\node[act] (mix) at (-4.2,-7.1) {Состав смеси\\или масса вещества};
\node[act] (price) at (4.2,-7.1) {Стоимость\\после изменения};
\node[act] (mixeq) at (-4.2,-9.5) {Таблица $m,\;c,\;q=mc$\\Баланс вещества};
\node[act] (priceeq) at (4.2,-9.5) {Таблица $S_0,\;p,\;S_1$\\Множитель $1\pm p/100$};
\node[act,text width=70mm] (solve) at (0,-11.8) {Составьте уравнение или систему\\и найдите неизвестное};
\node[act,text width=70mm] (check) at (0,-14.2) {Проверьте ограничения, единицы\\и ответьте на поставленный вопрос};
\node[start] (end) at (0,-15.9) {Запишите ответ};
\draw[line] (start)--(target);
\draw[line] (target)--(decision);
\draw[line] (decision.west)-| node[pos=0.22,above,font=\scriptsize] {смесь} (mix.north);
\draw[line] (decision.east)-| node[pos=0.22,above,font=\scriptsize] {цена} (price.north);
\draw[line] (mix)--(mixeq);
\draw[line] (price)--(priceeq);
\draw[line] (mixeq.south) |- (solve.west);
\draw[line] (priceeq.south) |- (solve.east);
\draw[line] (solve)--(check);
\draw[line] (check)--(end);
\end{tikzpicture}
\end{center}
\clearpage

\sectiontitle{Тренировочный блок --- Путь А}
\textcolor{softgray}{10 заданий по нарастающей сложности. Оформите уравнение или систему и запишите ответ с единицами измерения.}

\begin{enumerate}[label=\textbf{\arabic*.},leftmargin=8mm,itemsep=4mm]
\item После повышения цены на $12\%$ товар стал стоить $2240$ рублей. Какова была его первоначальная цена?
\item Цена товара снизилась с $3900$ до $3510$ рублей. На сколько процентов снизилась цена?
\item Смешали $12\unit{кг}$ $20\%$ раствора и $8\unit{кг}$ $45\%$ раствора одного вещества. Найдите концентрацию полученной смеси.
\item Первый сплав содержит $15\%$ меди, второй --- $35\%$. Масса второго на $4\unit{кг}$ больше массы первого, а полученный сплав содержит $30\%$ меди. Найдите массу полученного сплава.
\item Свежие яблоки содержат $80\%$ воды, сушёные --- $20\%$ воды. Сколько килограммов свежих яблок потребуется для получения $12\unit{кг}$ сушёных? Масса сухого вещества при сушке сохраняется.
\item Смешали $100\unit{кг}$ двух сплавов с содержанием меди $15\%$ и $35\%$. В новом сплаве $27\%$ меди. На сколько килограммов масса второго сплава больше массы первого?
\item Три одинаковые футболки стоят на $10\%$ дешевле рюкзака. На сколько процентов четыре такие футболки дороже рюкзака?
\item Цена товара сначала выросла на $10\%$, затем полученную цену уменьшили на $10\%$. Сколько стал стоить товар, если первоначально он стоил $2000$ рублей?
\item При смешивании равных масс двух растворов получается раствор с концентрацией $30\%$. Если взять одну часть первого раствора и три равные по массе части второго, получится раствор с концентрацией $35\%$. Сколько килограммов растворённого вещества содержится в $5\unit{кг}$ первого раствора?
\item Смешали $20\%$ и $50\%$ растворы кислоты и добавили $10\unit{кг}$ чистой воды. Получили раствор с концентрацией $28\%$. Если вместо воды к тем же исходным растворам добавить $10\unit{кг}$ $60\%$ раствора кислоты, концентрация полученной смеси составит $40\%$. Сколько килограммов $20\%$ раствора использовали?
\end{enumerate}
\clearpage

\sectiontitle{Ответы к тренировочному блоку}
\begin{center}
\begin{tabularx}{\linewidth}{@{}p{20mm}X@{}}
\toprule
№ & Ответ \\
\midrule
1 & $2000\unit{руб}$. \\
2 & $10\%$. \\
3 & $30\%$. \\
4 & $8\unit{кг}$. \\
5 & $48\unit{кг}$. \\
6 & $20\unit{кг}$. \\
7 & На $20\%$ дороже. \\
8 & $1980\unit{руб}$. \\
9 & $1\unit{кг}$ растворённого вещества. \\
10 & $20\unit{кг}$ $20\%$ раствора. \\
\bottomrule
\end{tabularx}
\end{center}

\vspace{6mm}
\key{Короткая самопроверка для сложных заданий}
\begin{itemize}
\item \textbf{№\,5.} Сухая масса: $0{,}20x=0{,}80\cdot12$.
\item \textbf{№\,6.} Если первого сплава $x\unit{кг}$, то второго $100-x\unit{кг}$; составьте баланс меди.
\item \textbf{№\,9.} Для долей $a,b$: $a+b=0{,}60$, $a+3b=1{,}40$; после нахождения $a$ вычислите $5a$.
\item \textbf{№\,10.} Для масс $x,y$ составьте два баланса:
\[
\begin{cases}
0{,}20x+0{,}50y=0{,}28(x+y+10),\\
0{,}20x+0{,}50y+6=0{,}40(x+y+10).
\end{cases}
\]
\end{itemize}
\end{document}
